\hypertarget{interpretation}{}
\hypertarget{how-to-read-the-papers-evidence-correctly}{%
\chapter{How to read the paper\textquotesingle s evidence
correctly}\label{how-to-read-the-papers-evidence-correctly}}

\begin{longtable}[]{@{}p{0.20\textwidth}p{0.34\textwidth}p{0.36\textwidth}@{}}
\toprule\noalign{}
Finding & What it supports & What it does not support \\
\midrule\noalign{}
\endhead
\bottomrule\noalign{}
\endlastfoot
Joint HAC/Wald p = 6.42e-10 & Not all FF3 alphas are zero & Positive
skill or stable ranking \\
Five raw alphas survive BH & Five model-relative intercept discoveries &
Five positive winners; four are negative \\
Two posterior intervals wholly positive & Two pooled alphas remain
positive under the EB model & Those portfolios are guaranteed to
outperform \\
Wide bootstrap rank intervals & Point ranks are uncertain & Scores are
useless; uncertainty is itself information \\
120-month rank persistence 0.156 & Weak decade-to-decade structural
persistence & Short-horizon 0.863 proves stability \\
Forward spread 243 bps/year & Past ordering has average future
information & Annual rolling timing is useful \\
Rolling increment vs expanding = \ensuremath{-}0.042 & Re-estimation adds no dynamic
predictive power & No stable cross-sectional information exists \\
FF3/FF5 rank correlation 0.875; max move 11 & Broad agreement plus
material benchmark risk & Alpha is benchmark-independent \\
One SFA FDR discovery & One residual series shows supported one-sided
asymmetry & A 25-portfolio efficiency league table \\
\end{longtable}

\textbf{The disciplined final statement.} There is cross-sectional FF3
alpha variation and a stable ordering with some forward information, but
ranks are uncertain, dynamic rolling re-estimation adds no value over an
expanding estimate, and conclusions depend on the benchmark and source
vintage.

